\section{Results}

In this section, I will present the results of the experiments previous conduced.
The main results are summarized in table \ref{tab:results_bench}.
The restults include the mean and standard deviation (SD) of the F1-Score, ROC AUC, and AUPRC for each method evaluated. 
To conclude my results I will apply a statistical test.
After presenting the results, I will display the confusion matrices for each method.
Additionally I will plot the the AUPRC and ROC AUC to further analyse the performance of each method.


%\begin{table}[h]
%\centering
%\caption{Experimental Results - Oneshot classification}
%\label{tab:results_os}
%\begin{tabular}{lcccccc} % l=left aligned, r=right aligned
%\toprule
%\textbf{Method}         & \textbf{F1-Score} &  \textbf{ROC AUC}         & \textbf{AUPRC}                        &        \\
%\midrule
%MCD                     & 50.00             & \textbf{97.56}            & \textbf{71.04}                        &        \\
%k-NN (unsupervised)     & 50.00             & 94.54                     & 50.51                                 &        \\
%\bottomrule
%\end{tabular}
%\end{table}

%The main finding is, that MCD performs better for this dataset that k-NN.
%The F1-Score is the same for both methods.
%The both deviate between the ROC AUC and the AUPRC. 
%ROC AUC is 3.02 percentage points bigger for MCD than k-NN, which is a slight difference.


%AUPRC is 20.53 percentage points bigger fpor MCD than k-NN which is a big difference and needs more investigation.
%That indicates, that even with the multidimensional skewed dataset, MCD is capable of determinang an unbiased 
\todo{sollte ich das wirklich mitreinnehmen? - Ja}


\begin{table}[h]
\centering
\caption{Experimental Results of benchmark protocol}
\label{tab:results_bench}
\begin{tabular}{lcccccc} 
\toprule
\textbf{Method}         & \textbf{F1-Score} & \textbf{SD} & \textbf{ROC AUC} & \textbf{SD} & \textbf{AUPRC} & \textbf{SD}          \\
\midrule
MCD  (semi-supervised)  & 64.14             & 9.41  &  95.94            & 3.59     & 77.82             & 16.43  \\
k-NN (semi-supervised)  & 68.12             & 8.62  &  95.87            & 1.89     & 72.85             & 11.54  \\
InterCont               & \textbf{77.44}    & 10.59 &  \textbf{96.70}   & 2.06     & \textbf{99.48}    & 3.2   \\
\bottomrule
\end{tabular}
\end{table}


The performance utilizes F1-Score, ROC AUC, and AUPRC, accompanied by their respective standard deviations (SD). 
In the F1-Score analysis, InterCont yields the highest mean at 77.44, followed by k-NN at 68.12 and MCD at 64.14. 
The SD values indicate that k-NN maintains the lowest variability at 8.62, while MCD and InterCont show higher fluctuations at 9.41 and 10.59.


The ROC AUC results show high discriminative performance across all methods, with InterCont reaching 96.70, k-NN at 95.87, and MCD at 95.94. 
The stability of these scores is evidenced by low SD values, with k-NN demonstrating the least variance at 1.89, compared to 2.06 for InterCont and 3.59 for MCD.


For the AUPRC metric, InterCont achieves a mean of 99.48, significantly higher than MCD at 77.82 and k-NN at 72.85. 
The InterCont method exhibits near-zero variance with an SD of 0.32. 
Conversely, MCD shows the highest instability in this metric with an SD of 16.43, while k-NN records an SD of 11.54. 
These figures reflect the comparative effectiveness and consistency of each method within the tested framework.



The authors of InterCont tested their results with a wilcoxon signed rank test. 
They conducted that by statistical sinificance, that their model outperforms other methods.
I can't apply a wilcoxon signed rank test to the results of my experiments, because this I applied only one dataset. \todo{argue why I dont chose Wilcoxon}


For the statistical analysis and to answer the research question , I decided to use a parametric approach.
Due to the central limit theorem and the number of runs (n=500) the distribution of the metrics will approximatly follow normal, regardless of the underlying distribution. \todo{really?}
Therefore I will deploy an One-Way ANOVA to analyse if there is a statistically significant difference between the methods.
If so, I will further investigate by apply Welchs's t-Test, to determine local superior performance.


The One-Way ANOVA, sometimes called ANOVA from summary statistics, tests if the means of the methods are differente, hence

\begin{equation}
    Ho = my1 = my2 = my3
\end{equation}

If the Null is rejected ($p<0.05$) that means that at least one mean diviates from the the other. 
It does not define which methods are different.
To estimate the F-statistics the  Between-Group Variance and Within-group variance has to be calculated.
The Between-Group variance measures the difference between the global mean and the methods mean of the specific metric.
This is the actual difference caused by the method.

The Within-Group Variance measures the variance inside the metric of a specific method.
Therefore I use the already estimated standard deviation. 


In the next step the mean squares of the between and within variance is estimated. 
Therefore I divide the variance by their degrees of freedom ($df$).
The $df$ is equal to the number of methods minus 1 for the Between-Group variance.
And for the Within-Group variance the $df$ is number of total runs - number of methods.
Because of the high number of runs, the denominator will be large, hence the mean squared will be small.
That means than even a small variance between the means will likely result in a high F value and a low p-value \todo{really}
The reliability of the SS within will be precise as well. \todo{why?}


The F-Value is calculated as follows:

\begin{equation}
    F = \frac{MS_between}{MS_within}
\end{equation}

That resembles the ratio of signal to noise.
If the F-Value is greater than $1$, the difference between the methods is larger than would be expected by random chance.
The results are displayed in \ref{tab:ANOVA}


\begin{table}[h]
\centering
\caption{Results of One-Way ANOVA}
\label{tab:ANOVA}
\begin{tabular}{lcccccc} 
\toprule
\textbf{Metric}         & F-Statistic       & P-Value   &   \\
\midrule
F1-Score                & 254.1747          & $9.12 \times 10^{-96}$     &   \\
ROC AUC                 & 15.3523           & $2.51 \times 10^{-07}$     &   \\
AUPRC                   & 727.5883          & $1.79 \times 10^{-221}$     &   \\
\bottomrule
\end{tabular}
\end{table}

To determine which specific methods performs for a specific metric better than the other, I will deploy a independent samples t-test, also known as Welchs's t-test.
The Welchs's t-test is a paranetric test to compare independend groups with different variances. 
Because I have 3 methods, I will perform 3 different pairs.
I chose this and not the well known stuendts t-test, because the students t-test assumes homoscedasticity.
The test assumes, that the runs are independently performed and the means are normally distributed.
Both assumptions are satisfied.
The test statistic is calculated as follows: 


\begin{equation}
    t = \frac{\bar{X}_1 - \bar{X}_2}{\sqrt{\frac{s_1^2}{n_1} + \frac{s_2^2}{n_2}}}
\end{equation}

The numerator resembles the signal, as in difference in means.
The denominator is the the noise, as in difference in estimazted standard error.
The $df$ are estimated, by adjusting the df based on the variance of each group to account for the fact that the two groups dont share a common variance.
Therefore the Welch-Satterthwaite equation estimates the number:

\begin{equation}
    \nu \approx \frac{\left( \frac{s_1^2}{n_1} + \frac{s_2^2}{n_2} \right)^2}{\frac{(s_1^2 / n_1)^2}{n_1 - 1} + \frac{(s_2^2 / n_2)^2}{n_2 - 1}}
\end{equation}

Therefore the $df$ can be decimal number.
It reflects how much effective information I have \todo{more precise}
When the variances of the two estmators are equal, the $df$ would be the same as in student's t-test.
With increasing difference in variance between two estimators, the lower the number of $df$ will be.
The result will be more conservative leading to a higher p-value.
Because $n=500$ the t-distirbution will follow a normal distribution.

Because I compare three different methods, I need to adjust the significance level.
The likelihood of finding a significant result would be increased, resultin in a potentially increased false positives. 
The Bonferroni corrected alpha is applied by simply deviding the significance level by the number of methods.
Therefore, to reject the Null, the p-value has to be lower than 0.0167. 
\todo{maybe insert familiy wise error rate (FWER)}


\begin{table}[h]
\centering
\caption{p-values of Welch's test for ROC AUC}
\label{tab:Welch1}
\begin{tabular}{lcccccc} 
\toprule
\textbf{ROC AUC}    & MCD       & k-NN      & InterCont  &    \\
\midrule
MCD                 & -             &               &    &     \\
k-NN                & 6.9975e-01    & -             &    &     \\
InterCont           & 4.4455e-05    & 5.2073e-11    & -  &     \\
\bottomrule
\end{tabular}
\end{table}


In the Table \ref{tab:Welch1} are the results of the Welches test for the ROC AUC.
It displays the only combination of methods, which result is not statistically different.
The ROC AUC mean of MCD and k-NN are statistically indistinguishable.
All other possible combinations are with probability of under 1.67\% different from each other.
The remaining result of the other metrics did not include exciting result and are included in the appendix. 
 
To conclude, the results of the InterCont were higher in every metric and can be interpreted as statistically different from the other.


\subsection{Discussion}
The goal is to write why the results fit into the theory.


The story is, that outlier detection should be implemented into realworld application.
Problem is the scarcity of labels. 
Therefore, 

I could not reproduce all results from the original paper.
For k-NN I only achieved similiar results, by using a semisupervised approach

Results that are worthy to discuss

1. MCD achieved the highest ROC AUC, but the deviation is not big
2. InterCont achieves the highest AUPRC, but MCD is better in ROC AUC.

Explanations:
MCD needs gaussian distribution feature 3 and 4 are nearly normally distributed, which could explain how the statistical fit likley contributed to high results.


Intercont success is down to handle multidimentional data through learning the manifold structure.
I could be because of the UMAP can handle the outlers successfully too and it is constructed in a simliar way.


k-NN can handle the data depending on the approach.
Unsupervised is worst perfomring method, but if you let the method get access to a bit of the information before estimation, the results increase very much.
The authors compared the semisupervised approach to the selfsupervised approach - is that a fair comparison? 
Why do I thinnk that this is not a fail comparison? - Because selfsupervised means that the algorithm creates its own "lables" by applying contrastive learning.
For semi-supervised approach you need labeled data wich in the real world is not immpossible but hard to get \todo{insert introduction/scarcity of Lables}


UMAP analysis showed that the outliers are tightly and distingly clustered and therefor the ROC AUC are so high. 
The feature Proline can be identified as a collective outlier. THe sliding window approach and the mahalanobis distance capture this therefore the results are better in MCD, InterCont then of k-NN.


InterCont shows the lowest SD because of the 500 random splits and the contrastive mechanism provided a stable learning signal. \todo{ist that comparable with semisupervised knn?}

\subsubsection*{F1 difference bweteen unsupervised and semi-supervised}
The difference in F1 score between unsupervised and semisupervised knn is die to the approach. 
By training only on normal data, semisupervised k-nn can establish a very stable and accurate reference for normal behavior. 
When the outliers are introduced, their distance to these pure neighbors is much more pronounced. 


By defining the outlier ratio to the actual outlier ratio, I force the models to get to the most optimal sport on the precision recall curve,
That maximized the F1-Score, becasue it is the harmonic mean of precision and recall (you can read about that on the authors website, where they published the paper)

\section{Limitations}

This chapter is about limitation of my thesis.
My approach is very dependend on the chosen datset.
In outlier detection, it is common practice to run the experiments on different datasets, to analyse the strengths and weaknesses of the method.
I did not sufficient hyperparamter tuning via phythons optuna.
I just made grid search which is not as concrete 


As a general detection performance result, we can conclude that nearest-neighbor based algorithms perform better in most cases when compared to clustering algorithms. 
Also, the stability concerning a not-perfect choice of k is much higher for the nearest-neighbor based methods. 
The reason for the higher variance in clustering-based algorithms is very likely due to the non-deterministic nature of the underlying k-means clustering algorithm. 
Despite of this disadvantage, clustering-based algorithms have a lower computation time. 
As a conclusion, we recommend to prefer nearest-neighbor based algorithms if computation time is not a issue. 
If a faster computation is required for large datasets, for example in a near real-time setting, clustering-based anomaly detection might be the method of choice. 
For small datasets, clustering based methods should be avoided.
Among the nearest-neighbor based methods, the global k-NN algorithm is a good candidate on average.


\todo{insert analysis of the outliers}
\todo{explain why MCD is bad for real world data, why k-NN is better for multimodel clusters. Why is contrastive learning the best choice for the real world dataset}
\todo{the resulting intercont scores are interpretable, k-NN is the distance interpretable}
